% shared by main.tex (IEEEtran) and access/main.tex (IEEE Access). Numbers come from generated/ only.
Can a language model small enough to live in the flash memory of a microcontroller write useful code? We study this question with an openly released, largely scripted setup: a benchmark of IoT-oriented Python function-synthesis tasks with executable unit tests and audited train/test separation, decoder-only Transformers with 0.53\,M and 1.05\,M parameters, and a portable INT8 C inference engine, verified on a PC and run inside a host simulation of the ESP32-S3 firmware (a JavaScript port for the browser reproduces its output on all test prompts). On unseen wordings of tasks from the training families, the 1.05\,M model passes the unit tests on \testidgreedypass\,\% of tasks with one-shot greedy decoding (mean of \nseeds{} seeds). \emph{Signature anchoring}, which forces the answer to start with the function header copied from the request, raises this to \testidanchorpass\,\%. A compiler-in-the-loop repair procedure (Recursive Error-Feedback, REF) gated by \texttt{compile()} and a static undefined-name check reaches \testidareflintpass\,\%. REF mainly fixes syntax: it is no better than resampling under the same budget, and replacing the error message with a constant placeholder does not measurably hurt. The 0.53\,M model is at least as accurate (\seedsSmalltestidanchor\,\% with anchoring, s.d.\ \seedsSmalltestidanchorsd) in \dsmflashtotalbytes{} bytes of flash. On task families held out from training, accuracy falls to \testoodgreedypass--\seedsSmalltestoodareflint\,\%, and on MBPP to zero. Million-parameter models are therefore practical, verifiable \emph{specialists} for a closed catalogue of embedded utility functions, not general programmers. We release all code, data, models and firmware.
